\documentclass[12pt]{article}
\usepackage[T1]{fontenc}
\usepackage[utf8]{inputenc}
\usepackage{lmodern}
\usepackage{textcomp}
\usepackage{enumitem}
\usepackage{geometry}
\geometry{margin=1in}
\begin{document}
\begin{center}
Answer Key - Economics - Undergraduate Level Assessment 17 \
\small (Answers are ordered exactly as in the question paper.)
\end{center}

\section*{Multiple Choice}
\begin{enumerate}[label=\arabic*.]
\item \textbf{Answer:} A
\\ \textit{Options:} A) Biased, inconsistent, and inefficient.; B) Biased but consistent; C) Unbiased and consistent; D) Biased but inconsistent for one equation only; E) Unbiased and efficient
\\ \textit{Note:} The correct answer is biased, inconsistent, and inefficient because applying Ordinary Least Squares (OLS) to each equation separately in a simultaneous equation system ignores the endogeneity caused by the simultaneous relationships among the variables, leading to inaccurate estimates.
\item \textbf{Answer:} A
\\ \textit{Options:} A) 5 percent \$475 in excess reserves; B) 10 percent \$50 in excess reserves; C) 30 percent \$350 in excess reserves; D) 60 percent \$200 in excess reserves; E) 80 percent \$100 in excess reserves
\\ \textit{Note:} The reserve ratio is calculated by dividing the required reserves ($50) by the total deposits ($500), resulting in 10 percent, and the excess reserves are found by subtracting the required reserves from the total deposits, yielding $450 in excess reserves, not $475 as stated in the answer.
\item \textbf{Answer:} B
\\ \textit{Options:} A) exports.; B) exports minus imports.; C) imports.; D) gross domestic product.; E) consumption.
\\ \textit{Note:} In the equation GDP = C + I + G + X, the term X represents net exports, which is calculated as exports minus imports, reflecting the value of goods and services produced domestically and sold abroad, adjusted for those purchased from foreign markets.
\item \textbf{Answer:} A
\\ \textit{Options:} A) decrease by \$200 million.; B) decrease by \$100 million.; C) increase by \$200 million.; D) decrease by \$10 million.; E) increase by \$500 million.
\\ \textit{Note:} The correct answer is based on the money multiplier effect, where a reserve requirement of five percent means that for every dollar the FED withdraws through the sale of securities, the total money supply can decrease by up to 20 times that amount, resulting in a potential decrease of $200 million from the initial $10 million sold.
\item \textbf{Answer:} A
\\ \textit{Options:} A) One can have comparative advantage in producing both goods.; B) One can have both an absolute advantage and a comparative advantage in producing x.; C) One can have absolute advantage and no comparative advantage in producing x.; D) One can have comparative advantage and no absolute advantage in producing x.
\\ \textit{Note:} The correct answer is that one can have comparative advantage in producing both goods because comparative advantage is based on the relative opportunity costs of production, and by definition, a producer can only have a comparative advantage in one good if they have a lower opportunity cost in that good compared to others.
\end{enumerate}

\section*{True / False}
\begin{enumerate}[label=\arabic*.]
\setcounter{enumi}{5}
\item \textbf{Answer:} False
\\ \textit{Options:} A) True; B) False
\\ \textit{Note:} The statement is false because a perfectly competitive firm's demand curve is horizontal at the market price, reflecting that the firm is a price taker and can sell as much as it wants at that price, rather than having a downward-sloping demand curve like the overall market demand.
\item \textbf{Answer:} False
\\ \textit{Options:} A) True; B) False
\\ \textit{Note:} This statement is false because the primary source of money creation in an economy is the banking system, specifically through the process of fractional reserve banking, where commercial banks create money by issuing loans.
\item \textbf{Answer:} True
\\ \textit{Options:} A) True; B) False
\\ \textit{Note:} The statement is true because an increase in demand raises the equilibrium quantity while the simultaneous increase in supply can offset the price increase from demand, leading to an uncertain effect on price.
\item \textbf{Answer:} True
\\ \textit{Options:} A) True; B) False
\\ \textit{Note:} Conciliation, mediation, voluntary arbitration, and compulsory arbitration are recognized as effective methods for resolving labor-management impasses because they provide structured processes that facilitate negotiation, promote communication, and offer binding or non-binding resolutions that help to prevent prolonged disputes and strikes.
\item \textbf{Answer:} True
\\ \textit{Options:} A) True; B) False
\\ \textit{Note:} The statement is true because, when adjusting for 12\% annual inflation over 4 years, the purchasing power of cash balances is reduced, and the calculation shows that the real value would be approximately \$599.69.
\end{enumerate}

\section*{Long Form Response}
\begin{enumerate}[label=\arabic*.]
\setcounter{enumi}{10}
\item \textbf{Answer:} Currency exchange rates reflect the relative value of one currency to another, and discrepancies in these rates can create arbitrage opportunities. For instance, if the price of a British pound is $1.00 and the price of a German mark is $0.40, the expected exchange rate would be 2.5 marks per pound (1.00 / 0.40). If the actual market exchange rate deviates from this expected rate, an Englishman could profit by buying pounds in the cheaper market and selling them in the more expensive market. This profit arises from exploiting the differences in exchange rates to realize gains without risk, highlighting the dynamic nature of currency markets.
\\ \textit{Note:} The answer correctly identifies that discrepancies in currency exchange rates create arbitrage opportunities by allowing individuals to profit from buying low in one market and selling high in another, supported by the calculation of the expected exchange rate based on the given prices.
\item \textbf{Answer:} Keynesian economics fundamentally supports a market economy, positing that while private enterprise plays a crucial role in economic growth, government intervention is necessary to stabilize and guide the economy during periods of recession or inflation. Keynes did not oppose market principles; rather, he argued that active government involvement can complement market forces by addressing failures and ensuring full employment. This dual approach suggests that both government and private enterprise are essential for a balanced and efficient economy. By advocating for government intervention, Keynes aimed to enhance market functionality rather than undermine it, asserting that a collaborative effort between the two can lead to greater economic stability and prosperity.
\\ \textit{Note:} The answer correctly highlights that Keynesian economics advocates for a balanced coexistence of market forces and government intervention, emphasizing that such intervention is designed to stabilize the economy without opposing the fundamental principles of a market economy.
\end{enumerate}

\vfill
\noindent\textit{End of Answer Key}
\end{document}